\begin{center}
{\Large \bfseries \textcolor{lyred}{第二章 导数与微分}}
\end{center}

\vspace{6pt}
{\large \bfseries \textcolor{lyred}{第2.1 节 导数概念}}

\vspace{4pt}
{\bfseries \textcolor{lyred}{一.引例}}

{\bfseries \textcolor{lyred}{1.直线运动的瞬时速度}}

{\bfseries \textcolor{lyred}{2.平面曲线的切线斜率}}

\vspace{4pt}
{\bfseries \textcolor{lyred}{二.导数的定义}}

{\bfseries \textcolor{lyred}{1.函数在一点处的导数与导函数}}

\textcolor{lyred}{定义}.设
()
y
f x
=
在某个
(
)
U x
内有定义,称它在
0x 处\textcolor{lyred}{可导},若极限 
(
)
(
)
()
(
)
lim
lim
lim
x
x
x
x
f x
x
f x
f x
f x
y
x
x
x
x
\Delta \to 
\Delta \to 
\to 
+\Delta 
-
-
\Delta =
=
\Delta 
\Delta 
-
存在,称它为
()
f x 在
0x 处的 
\textcolor{lyred}{导数},记为
(
)
f
x
'
,它反映了
()
f x 在
0x 处随x 变化的快慢,故也称为\textcolor{lyred}{变化率},它是 
平均变化率的极限. 
\textcolor{lyred}{定义}.若
()
f x 在(
)
,a b 内处处\textcolor{lyred}{可导},则称它为(
)
,a b 上的\textcolor{lyred}{可导函数},记 
()
(
)
()
lim
x
f x
x
f x
f
x
x
\Delta \to 
+\Delta 
-
'
=
\Delta 
,称为\textcolor{lyred}{导函数},简称\textcolor{lyred}{导数},也记为y', df
dx , dy
dx . 
{\bfseries \textcolor{lyred}{2.求导举例}}

\textcolor{lyred}{例}.()
lim
h
C
C
C
h
\to 
-
'=
=
,()
(
)
lim
h
x
h
x
x
h
\to 
+
-
'=
=. 
\textcolor{lyred}{例}.(
)
(
)
(
)
lim
lim
lim 2
h
h
h
x
h
x
xh
h
x
x
h
x
h
h
\to 
\to 
\to 
+
-
+
'=
=
=
+
=
. 
\textcolor{lyred}{例}.(
)
(
)
lim
lim
lim
h
h
h
h
h
x
h
x
x
x
x
x
x
x
h
h
h
\mu 
\mu 
\mu 
\mu 
\mu 
\mu 
\mu 
\mu 
\mu 
-
\to 
\to 
\to 


+
-
\cdots 


+
-
'


=
=
=
=
. 
\textcolor{lyred}{注}.特别地,(
)
lim
h
x
h
x
x
h
x
\to 
+
-
'=
=
,
lim
h
x
h
x
x
h
x
\to 
-
'


+
=
=-




. 
\textcolor{lyred}{例}.(
)
(
)
2cos
sin
sin
sin
sin
lim
lim
cos
h
h
h
h
x
x
h
x
x
x
h
h
\to 
\to 


+


+
-


'=
=
=
; 
类似地,(
)
cos
sin
x
x
'=-
. 
\textcolor{lyred}{例}.(
)
lim
lim
ln
x h
x
h
x
x
x
h
h
a
a
a
a
a
a
a
h
h
+
\to 
\to 
-
-
'=
=
=
,于是(
)
x
x
e
e
'=
. 
\textcolor{lyred}{例}.(
)
(
)
log
ln 1
log
log
log
lim
lim
lim
ln
a
a
a
a
h
h
h
h
h
x
h
x
x
x
x
h
h
h
a
\to 
\to 
\to 




+
+




+
-




'=
=
=
= 
lim
ln
ln
h
h
x
h
a
x
a
\to 
=
,于是(
)
ln x
x
'=
. 
{\bfseries \textcolor{lyred}{3.导数的运动学意义}}

设
()
s
s t
=
为直线运动物体的位置函数,则
ds
v
dt
=
为\textcolor{lyred}{速度}(位置对时间的变化率), 
反映物体运动的快慢.
dv
a
dt
=
为\textcolor{lyred}{加速度},反映物体运动速度变化的快慢. 
\vspace{4pt}
{\bfseries \textcolor{lyred}{三.导数的几何意义}}

平面曲线C :
()
y
f x
=
在
(
)
,
M x
y
处的切线斜率
(
)
()
(
)
lim
x
x
f x
f x
f
x
x
x
\to 
-
'
=
-
,故 
切线方程为
(
)(
)
(
)
y
f
x
x
x
f x
'
=
-
+
,法线方程为
(
)(
)
(
)
y
x
x
f x
f
x
-
=
-
+
'
. 
\textcolor{lyred}{例}.求曲线
y
x
=
上与直线
y
x
=
-平行的切线方程. 
解.设切点为(
)
,
x
y
,则
x
x
y
=

=
\Rightarrow 
=

,故切线方程为
(
)
y
x
=
-
+. 
\textcolor{lyred}{例}.求曲线
y
x
=
上通过(
)
0, 4
-
的切线方程. 
解.设切点为
,
x x






,则
(
)
x
x
x
x
--
=
\Rightarrow 
=
-
,故切线方程为
(
)
y
x
=
-
+. 
\vspace{4pt}
{\bfseries \textcolor{lyred}{四.函数可导性与连续性的关系}}

\textcolor{lyred}{定理}.设
()
f x 在
0x 处可导,则它在
0x 处连续. 
\textcolor{lyred}{证}.
()
(
)
()
(
)(
)
(
)
lim
lim
x
x
x
x
f x
f x
f x
f x
x
x
f
x
x
x
\to 
\to 
-
'
-
=
-
=
\cdots 
=




-
,证毕. 
\textcolor{lyred}{注}.反之不对,例如
()
f x
x
x
=
-
在
0x 处连续,但是在
0x 处不可导. 
\textcolor{lyred}{例}.设
()
sin
,
0,
nx
x
f x
x
x

\neq 

=

=

,其中n 为正整数.讨论 
(1)
()
f x 在
x =
处的连续性;(2)
()
f x 在
x =
处的可导性. 
解.(1)
()
lim
lim
sin
n
x
x
f x
x
x
\to 
\to 
=
=
,故连续; 
(2)
()
()
lim
lim
sin
n
x
x
f x
f
x
x
x
-
\to 
\to 
-
=
-
,故当
n >时可导,当
n =时不可导. 
\vspace{4pt}
{\bfseries \textcolor{lyred}{五.单侧导数}}

\textcolor{lyred}{定义}.记
(
)
()
(
)
lim
x
x
f x
f x
f
x
x
x
-
-
\to 
-
'
=
-
,称为
()
f x 在
0x 处的\textcolor{lyred}{左导数}; 
记
(
)
()
(
)
lim
x
x
f x
f x
f
x
x
x
+
+
\to 
-
'
=
-
,称为
()
f x 在
0x 处的\textcolor{lyred}{右导数}. 
\textcolor{lyred}{定理}.
()
f x 在
0x 处可导
(
)
(
)
f
x
f
x
-
+
'
'
\Leftrightarrow 
=
. 
\textcolor{lyred}{例}.设
()
3,
2,
x
x
f x
x
x

\le 
=
-
>

,讨论
()
f x 在
x =处的可导性. 
解.
()
lim
x
x
f
x
-
-
\to 
-
'
=
=
-
,
()
(
)
lim
x
x
f
x
+
+
\to 
-
-
'
=
=
-
,故可导. 
\textcolor{lyred}{例}.设
()
(
)
,
ln
,
x
a
x
f x
x
e
b
x

\le 

=
+
+
>

,求a ,b ,使得
()
f x 在
x =
处可导. 
解.可导\Rightarrow 连续,由(
)
(
)
()
f
f
f
-
+
=
=
,得
b =
; 
()
lim
ln
x
x
a
f
a
x
-
-
\to 
-
'
=
=
-
,
()
(
)
ln
lim
lim
ln
x
x
x
e
x
f
x
x
e
e
+
+
+
\to 
\to 
+
-


'
=
=
+
=


-


,得
e
a
e
=
. 
\textcolor{lyred}{定义}.设
()
f x 在(
)
,a b 内可导,并且
()
f
a
+'
与
()
f
b
-'
均存在,则称
()
f x 为[
]
,a b 上的 
\textcolor{lyred}{可导函数}.此时,它一定在[
]
,a b 上连续. 
\vspace{4pt}
{\bfseries \textcolor{lyred}{六.例题}}

\textcolor{lyred}{例}.设
()
(
)
lim
x
f
f
x
x
\to 
-
-
=-,求
()1
f '
. 
解.
()
(
)
(
)
()
()
lim
lim
x
x
f
f
x
f
x
f
f
x
x
\to 
\to 
-
-
-
-
'
=
=
-
,故
()1
f '
=-. 
\textcolor{lyred}{例}.设
(
)
f
x
A
'
=
,求
(
)
(
)
lim
h
f x
h
f x
h
h
\to 
+
-
-
. 
解.上式
(
)
(
)
(
)
(
)
(
)
1 lim
2 h
f x
h
f x
f x
f x
h
f
x
A
h
h
\to 
+
-
-
-


'
=
+
=
=




. 
\textcolor{lyred}{例}.设
()
f x
x
x
=
-
,则
(
)
(
)
lim
lim
h
h
h
h
f x
h
f x
h
h
h
\to 
\to 
--
+
-
-
=
=
,但是 
(
)
(
)
(
)
lim
lim
h
h
h
f x
h
f x
f
x
h
h
\to 
\to 
+
-
'
=
=
,不存在. 
\textcolor{lyred}{例}.设
()
f x 在
x =
处连续,且当
x \to 
时,
()
()
f x
o x
=
,求
()
f '
. 
解.
()
()
()
lim
x
f x
f x
o x
x
\to 
=
\Rightarrow 
=
,故
()
()
lim
x
f x
f
x
\to 
'
=
=
. 
\textcolor{lyred}{例}.设
()
f x 在
x =
处连续,
()
lim
x
f x
x
\to 
=
-
,求
()
f '
. 
解.
()
()
()
1 lim
lim
x
x
f x
f x
f
x
x
\to 
\to 
=\Rightarrow 
=
\Rightarrow 
=
-
-
,
()
()
lim
x
f x
f
x
\to 
'
=
=
-
. 
\textcolor{lyred}{注}.若
()
f x 在
0x 处连续,
()
lim
x
x
f x
A
x
x
\to 
=
-
,则
(
)
()
lim
x
x
f x
f x
\to 
=
=
,
(
)
f
x
A
'
=
. 
\textcolor{lyred}{例}.设曲线
()
y
f x
=
在原点处与正弦曲线相切,求
lim
n
nf
n
\to \infty 






. 
解.
()
f
=
,
()
f '
=,故
()
lim
x
f x
x
\to 
=,
()
lim
2lim
n
x
f x
nf
n
x
\to \infty 
\to 

=
=




. 
\textcolor{lyred}{例}.设
()
f x 在x
a
=
处可导,且
()
f a =
,令
()
()
F x
f x
=
,证明:
()
F x 在x
a
=
处 
可导
()
f
a
'
\Leftrightarrow 
=
. 
证.
()
()
()
()
lim
lim
x
a
x
a
f x
f x
F
a
f
a
x
a
x
a
-
-
-
\to 
\to 
'
'
=
=-
=-
-
-
, 
()
()
()
()
+
lim
lim
x
a
x
a
f x
f x
F
a
f
a
x
a
x
a
+
+
\to 
\to 
'
'
=
=
=
-
-
,证毕. 
\textcolor{lyred}{补充练习 }
1.设
()
()
cos
,
0,
g x
x
f x
x
x

\neq 

=

=

,且()
()
g
g'
=
=
,求
()
f '
. 
解.
()
lim
x
g x
x
\to 
=
,故
()
()
()
lim
lim
cos
x
x
f x
g x
f
x
x
x
\to 
\to 
'
=
=
=
. 
2.设()
x
\varphi 
在x
a
=
处连续,令
()
()
f x
x
a
x
\varphi 
=
-
,证明:
()
f x 在x
a
=
处可导\Leftrightarrow  
()
a
\varphi 
=
. 
证.
()
()
()
()
lim
lim
x
a
x
a
x
a
f x
f a
x
a
x
a
x
a \varphi 
\varphi 
\to 
\to 
-
-
=
=\pm 
-
-
,即得,证毕. 
3.设
()
f '
=
,求
(
)
(
)
(
)
sin
tan
lim
ln 1 2
x
f
x
f
x
x
\to 
-
-
+
. 
解.原式
(
)
()
(
)
()
()
()
sin
tan
lim
x
f
x
f
f
x
f
f
f
x
\to 
-
-
-
-


'
'


=
=
+
=
. 
4.设
()
f x 在
x
x
=
处可导,求
(
)
()
lim
x
x
xf x
x f x
x
x
\to 
-
-
. 
解.原式
(
)
(
)
(
)
()
(
)
(
)
lim
x
x
xf x
x f x
x f x
x f x
f x
x f
x
x
x
\to 
-
+
-
'
=
=
-
-
. 
5.设
()
f x 在
x =
处连续,
(
)
lim
2sin
x
f
x
x
\to 
-
-
=
,求
()
f '
. 
解.由连续性,
()
(
)
lim
x
f
f
x
\to 
=
-
=,于是
(
)
()
lim
x
f
x
f
x
\to 
-
-
=
= 
(
)
()
()
lim
x
f
x
f
f
x
\to 
-
-
'
-
=-
-
,故
()
f '
=-
. 
6.设
()
,
arctan
0,
x
e
x
f x
x
x

-
\neq 
+
=

=

,讨论
()
f x 在
x =
处的可导性. 
解.
()
()
lim
x
f x
f
x
\pi 
-
-
\to 
'
=
=
-
,
()
()
lim
x
f x
f
x
\pi 
+
+
\to 
'
=
=
+
,故不可导. 
7.设
()
f
=,
()
f '
=
,求
()
1 cos
lim
x
x
x
f x
-
\to 



. 
解.
()
()
()
()
()
()
ln
2lim
2lim
1 cos
1 cos
lim
lim
x
x
x
f x
f x
f x
f
x
f
x
x
x
x
x
x
f x
e
e
e
e
e
\to 
\to 
-
-
'
-
-
-
\to 
\to 
=
=
=
=
=




. 
8.设
()
f x 连续,
()
lim
x
xf x
x
x
\to 
+
-
-=
,求
()
f '
. 
解.
()
()
()
()
1 2
1 2
xf x
x
x
o
f x
x
o x
x
x
+
-
-
-
-
=
+
\Rightarrow 
=
+
+
,故 
()
()
lim
x
f
f x
\to 
=
=,
()
()
1 2
lim
lim
x
x
f x
x
x
f
x
x
\to 
\to 
-
-
-
-
'
=
=
+
=. 
9.设
()
f x 在[
]
,a b 上可导,
()
()
f
a
f
b
+
-
'
'
\cdots 
<
,证明:
(
)
,a b
\psi 
\exists \in 
,使得
()
f
\psi 
'
=
. 
证.不妨设
()
()
()
lim
x
a
f x
f a
f
a
x
a
+
+
\to 
-
'
=
>
-
,
()
()
()
lim
x
b
f x
f b
f
b
x
b
-
-
\to 
-
'
=
<
-
,则 
在x
a
=
的右侧邻域
()
()
f x
f a
>
,在x
b
=
的左侧邻域
()
()
f x
f b
>
,故 
存在
(
)
,
x
a b
\in 
,使得
(
)
()
max
a x b
f x
f x
\le \le 
=
; 
由于
(
)
()
(
)
lim
x
x
f x
f x
f
x
x
x
+
+
\to 
-
'
=
\le 
-
,
(
)
()
(
)
lim
x
x
f x
f x
f
x
x
x
-
-
\to 
-
'
=
\ge 
-
,故 
(
)
(
)
(
)
f
x
f
x
f
x
+
-
'
'
'
=
=
=
,证毕. 
\vspace{6pt}
{\large \bfseries \textcolor{lyred}{第2.2 节 函数的求导法则}}

\vspace{4pt}
{\bfseries \textcolor{lyred}{一.函数和差积商的求导法则}}

\textcolor{lyred}{定理}.设()
u x 及()
v x 在
0x 处可导,则它们的和差积商(分母非零)也在
0x 处可导, 
并且(1)(
)(
)
(
)
(
)
u
v
x
u
x
v x
'
'
'
\pm 
=
\pm 
; 
(2)(
)(
)
(
)(
)
(
)(
)
uv
x
u
x
v x
u x
v x
'
'
'
=
+
; 
(3)
(
)
(
)(
)
(
)(
)
(
)
u
x
v x
u x
v x
u
x
v
v
x
'
'
'
-


=




,特别地,
(
)
(
)
(
)
v x
x
v
v
x
'
'
-


=




. 
\textcolor{lyred}{推论}.设()
u x , ()
v x 均在I 上可导,则在I 上,有(1)(
)
u
v
u
v
\alpha 
\beta 
\alpha 
\beta 
'
'
'
\pm 
=
\pm 
; 
(2)(
)
uv
u v
uv
'
'
'
=
+
,(
)
uvw
u vw
uv w
uvw
'
'
'
'
=
+
+
;(3)
u
u v
uv
v
v
'
'
'
-

=




. 
\textcolor{lyred}{例}.
4cos
sin 3
y
x
x
\pi 
=
+
-
,
(
)
(
)
4 cos
sin
4sin
y
x
x
x
x
\pi '


'
'
'=
+
-
=
-




. 
\textcolor{lyred}{例}.(
)
()
(
)
(
)
(
)
ln
ln
ln
ln
ln
ln
x
x
x
x
x
xe
x
x
e
x
x e
x
xe
x
e
x
x
x
'
'
'
'
=
+
+
=
+
+
. 
\textcolor{lyred}{例}.(
)
(
)
sin
sec
sec tan
cos
cos
x
x
x
x
x
x
'
--


'=
=
=




,(
)
csc
csc cot
x
x
x
'=-
. 
\textcolor{lyred}{例}.(
)
sin
cos
cos
sin
sin
tan
sec
cos
cos
x
x
x
x
x
x
x
x
x
'
\cdots 
+
\cdots 


'=
=
=




,(
)
cot
csc
x
x
'=-
. 
\textcolor{lyred}{例}.设
()
(
)()
f x
x
a
x
\varphi 
=
-
,其中()
x
\varphi 
在x
a
=
处连续,求
()
f
a
'
. 
解.
()
()
()
(
)()
()
()
lim
lim
lim
x
a
x
a
x
a
f x
f a
x
a
x
f
a
x
a
x
a
x
a
\varphi 
\varphi 
\varphi 
\to 
\to 
\to 
-
-
'
=
=
=
=
-
-
. 
\vspace{4pt}
{\bfseries \textcolor{lyred}{二.反函数的求导法则}}

\textcolor{lyred}{定理}.设
()
x
f
y
=
单调可导,且
()
f
y
'
\neq 
,则它的反函数
()
y
f
x
-
=
也可导,并且 
()
()
()
f
x
f
y
f
f
x
-
-
'

=
=


'
'



,即
dy
dx
dx
dy
=
. 
\textcolor{lyred}{例}.设
()
y
f x
=
为
x
y
y
y
=
+
+
+
的反函数,求
()
f '
. 
解.当
x =
时,
y =
,故
x
y
dy
dx
y
y
=
=
=
=
+
+
. 
\textcolor{lyred}{例}.
arcsin
y
x
=
,(
)
(
)
arcsin
cos
1 sin
sin
x
y
y
x
y
'=
=
=
=
'
-
-
; 
类似地,(
)
arccos
x
x
'=-
-
,或利用arccos
arcsin
x
x
\pi 
=
-
也可得. 
\textcolor{lyred}{例}.
arctan
y
x
=
,(
)
(
)
arctan
sec
tan
tan
x
y
y
x
y
'=
=
=
=
+
+
'
; 
类似地,(
)
arccot
x
x
'=-+
,或利用arccot
arctan
x
x
\pi 
=
-
也可得. 
\vspace{4pt}
{\bfseries \textcolor{lyred}{三.基本初等函数的导数公式}}

()
C '=
,(
)
x
x
\mu 
\mu 
\mu 
-
'=
,(
)
sin
cos
x
x
'=
,(
)
cos
sin
x
x
'=-
,(
)
tan
sec
x
x
'=
, 
(
)
cot
csc
x
x
'=-
,(
)
sec
sec tan
x
x
x
'=
,(
)
csc
csc cot
x
x
x
'=-
,(
)
x
x
e
e
'=
, 
(
)
ln
x
x
a
a
a
'=
,(
)
ln x
x
'=
,(
)
log
ln
a x
x
a
'=
,(
)
arcsin
x
x
'=
-
, 
(
)
arccos
x
x
'=-
-
,(
)
arctan
x
x
'=+
,(
)
arccot
x
x
'=-+
. 
\vspace{4pt}
{\bfseries \textcolor{lyred}{四.复合函数的求导法则(链式法则)}}

\textcolor{lyred}{定理}.设
()
u
x
\varphi 
=
在
0x 处可导,而
()
y
f u
=
在
(
)
u
x
\varphi 
=
处可导,则
()
y
f
x
\varphi 
=



 
在
0x 处可导,且
(
)
(
)
x
u
x
dy
dy
du
f
x
x
dx
du
dx
\varphi 
\varphi 
'
'
=
\cdots 
=




. 
\textcolor{lyred}{推论}.设可导函数
()
y
f u
=
,
()
u
x
\varphi 
=
可以复合,则
()
y
f
x
\varphi 
=



可导,并且 
()
()
()
()
dy
dy du
f
u
x
f
x
x
dx
du dx
\varphi 
\varphi 
\varphi 
'
'
'
'
=
\cdots 
=
\cdots 
=




. 
\textcolor{lyred}{注}.设
()
y
f u
=
,
()
u
g v
=
,
()
v
h x
=
可以复合,则dy
dy du dv
dx
du dv dx
=
\cdots 
\cdots 
= 
()
(
)
(
)
()
(
)
()
f
g h x
g h x
h x
'
'
'
\cdots 
\cdots 
. 
\textcolor{lyred}{注}.
()
()
()
f
x
f x
f
x
'
'

=
\cdots 


,
()
()
()
f
x
f x
f x
'
'

=


,
()
()
()
f
x
f x
f
x
'


'
-
=




, 
()
()
()
sin
cos
f x
f x
f
x
'
'
=
\cdots 




,
()
()
()
f x
f x
e
f
x e
'


'
=


,
()
()
()
ln
f
x
f x
f x
'
'=




, 
()
()
()
arctan
f
x
f x
f
x
'
'=




+
,等等. 
\textcolor{lyred}{例}.
tan
y
x
=
:
y
u
=
,
tan
u
x
=
,
3tan
sec
y
x
x
'=
\cdots 
. 
\textcolor{lyred}{例}.
arcsin
y
x
=
:
arcsin
y
u
=
,
u
x
=
,
x
y
x
'=
-
. 
\textcolor{lyred}{例}.
sin
x
x
y
e
-
=
:
u
y
e
=
,u
v
=
,
sin
v
x
x
=
-
,
sin
1 cos
sin
x
x
x
y
e
x
x
-
-
'=
\cdots 
-
. 
\textcolor{lyred}{例}.
y
x
a
=
-
,
(
)
3 2
x
x
y
x
a
x
a
x
a
-
-
'=
\cdots 
=
-
-
-
. 
\textcolor{lyred}{例}.
(
)
ln
y
x
x
=
+
+
,
(
)
ln
x
x
x
x
x
x
x
'




+
+
=
\cdots 
+
=






+
+
+
+


. 
\textcolor{lyred}{例}.
x
a
a
a
x
a
y
a
a
x
=
+
+
,
(
)
(
)
ln
ln
x
a
a
a
x
x
a
a
a
y
a
a
a
a
a
x
a x
-
'
'
'=
\cdots 
+
\cdots 
+
= 
ln
ln
x
a
a
a
x
x
a
a
a
a
a
a
x
a
a x
+
+
-
-
+
+
. 
\textcolor{lyred}{例}.
sin x
y
x
=
,
(
)
(
)
sin ln
sin ln
sin
sin
sin ln
cos ln
x
x
x
x
x
x
y
e
e
x
x
x
x
x
x


'
'
'=
=
\cdots 
=
+




. 
\textcolor{lyred}{例}.
ln ln ln tan 2
x
y




=








,
sec 2 2
ln ln tan
ln tan
tan
x
y
x
x
x
'=
\cdots 
\cdots 
\cdots 
\cdots 
. 
\textcolor{lyred}{例}.
ln cos
x
y
e
=
,
sin
2ln cos
tan
ln cos
cos
x
x
x
x
x
x
x
e
y
e
e
e
e
e
e
-
'=
\cdots 
\cdots 
\cdots 
=-
\cdots 
. 
\vspace{4pt}
{\bfseries \textcolor{lyred}{五.双曲函数的导数}}

(
)
sh
ch
x
x
'=
;(
)
ch
sh
x
x
'=
;(
)
th
ch
x
x
'=
; 
(
)
arsh
x
x
'=
+
;(
)
arch
x
x
'=
-
;(
)
arth
x
x
'=-
. 
\textcolor{lyred}{补充练习 }
1.设
()
sin 2
f x
x
=
,求
()
f
f x
'


,
()
\{
\}
f
f x
'




. 
解.
()
(
)
2cos 2sin 2
f
f x
x
'
=




,
()
\{
\}
(
)
2cos 2sin 2
2cos2
f
f x
x
x
'=
\cdots 




. 
2.设
(
)
ln tan
x
d
f
x
dx
\pi 
=
=




,求
()
f '
. 
解.
(
)
(
)
()
sec
ln tan
ln tan
tan
x
x
d
x
f
x
f
x
f
dx
x
\pi 
\pi 
=
=
'
'
=
\cdots 
=
=




,故
()
f '
=
. 
3.设
()
tan
tan
tan
x
x
x
f x
\pi 
\pi 
\pi 






=
-
-
-











\Lambda 
,求
()
f '
. 
解.
()
()
tan
x
f x
x
\pi 
\varphi 


=
-




,故
()
()
()
sec
tan
x
x
f
x
x
x
\pi 
\pi 
\pi 
\varphi 
\varphi 


'
'
=
\cdots 
+
-




, 
()
()
sec
99!
f
\pi 
\pi 
\pi 
\varphi 
'
=
\cdots 
=-
\cdots 
. 
4.设
()
f x 在
x =
处可导,且
()
f
=
,
()
f '
=-,求
()
lim
x
f
x
x
\to 
-
. 
解.
()
()
()
()
()
lim
lim
x
x
f
x
f
x
f
f
f
x
x
\to 
\to 
-
-
'
=
=
=-
-
. 
\vspace{6pt}
{\large \bfseries \textcolor{lyred}{第2.3 节 高阶导数}}

若
()
y
f x
=
的导数
()
f
x
'
仍可导,则它的导数称为
()
f x 的\textcolor{lyred}{二阶导数},记为
()
f
x
''
, 
y '',
d f
dx
,
dx
y
d
;二阶导数的导数称为\textcolor{lyred}{三阶导数},记为
()
f
x
'''
, y''',
d f
dx ,
d y
dx ; 
一般地,
n -阶导数的导数称为\textcolor{lyred}{n 阶导数},记为
()()
n
f
x ,
()
n
y
,
n
n
d f
dx
,
n
n
d y
dx . 
\textcolor{lyred}{法则1(线性性)}.(
)
()
()
()
()
1 1
1 1
n
n
n
n
m
m
m
m
u
u
u
u
u
u
\alpha 
\alpha 
\alpha 
\alpha 
\alpha 
\alpha 
+
+
+
=
+
+
+
\Lambda 
\Lambda 
. 
\textcolor{lyred}{法则2(莱布尼茨)}.(
)
()
()
(
)
(
)
()
n
n
n
n
n
n
n
n
uv
uv
u v
u v
u
v
n
-
-



'
''
=
+
+
+
+
=






\Lambda 
()(
)
n
k
n k
k
n u
v
k
-
=



\sum 
,其中
(
)
(
)
(
)
!
!
! !
n
n n
n
k
n
k
k
n
k
k
-
-
+
=
=

-

\Lambda 
. 
\textcolor{lyred}{例}.设
2x
y
x e
=
,求
(
)
y
. 
解.
(
)
(
)
(
)
(
)
(
)
(
)(
)
(
)
(
)(
)
(
)
20 19
2!
x
x
x
x
y
x
e
x
e
x
e
x
e
\cdots 
'
''
=
\cdots 
=
\cdots 
+
+
= 
(
)
20 2
190 2
x
x
x
x
x e
xe
e
x
x
e
+
\cdots 
\cdots 
+
\cdots 
\cdots 
=
+
+
. 
\textcolor{lyred}{例}.设
(
)
n
m
y
x
=
-
,求
()()1
n
y
. 
解.令
(
)(
)
()()
n
n
m
y
x
x
x
x
u x v x
-
=
-
+
+
+
+
=
\Lambda 
,由于
()()
k
u
=
,0
k
n
\le 
\le 
-, 
故
()()
()()
(
)()
()()()
(
)
!
n
n
n
k
n k
n
n
k
n
y
u
v
u
v
n m
k
-
=

=
=
=-
\cdots 


\sum 
. 
\textcolor{lyred}{法则3}.设
()
y
y x
=
二阶可导,且
()
y x
'
\neq 
,则
d x
y
dy
y
''
=-
'. 
\textcolor{lyred}{例}.设
x
y
sin
=
,求
()
n
y
. 
解.
cos
sin
y
x
x
\pi 


'=
=
+




,
sin
y
x
\pi 
\pi 


''=
+
+




,故
()
sin
n
y
x
n \pi 


=
+
\cdots 




. 
\textcolor{lyred}{注}.(
)
()
cos
cos
n
x
x
n \pi 


=
+
\cdots 




. 
\textcolor{lyred}{例}.设
sin
cos
y
x
x
=
+
,求
()
n
y
. 
解.
(
)
sin
cos
2sin
cos
sin 2
cos4
y
x
x
x
x
x
x
=
+
-
=-
=
+
,故 
cos 4
y
x
\pi 


'=
+




,
4cos 4
2 2
y
x
\pi 


''=
+\cdots 




,
()
cos 4
n
n
n
y
x
\pi 
-


=
+




. 
\textcolor{lyred}{例}.设
y
x
a
=
+
,求
()
n
y
. 
解.
(
)(
)
y
x
a
-
'=-
+
,
(
)(
)(
)
y
x
a
-
''=-
-
+
,
()
(
)
(
)
!
n
n
n
n
y
x
a
+
-
=
+
. 
\textcolor{lyred}{例}.设
y
x
x
=
-
+
,求
()
n
y
. 
解.
()
()
()
()
(
)
(
)
(
)
(
)
!
!
n
n
n
n
n
n
n
n
n
n
y
x
x
x
x
x
x
+
+
-
-






=
-
=
-
=
-






-
-
-
-






-
-
. 
\textcolor{lyred}{注}.
x
x
x
x
x
+
=
-
-
+
-
-
. 
\textcolor{lyred}{例}.设
()
()
f
x
f
x
'
=
,
()
f
=,求
()()
n
f
. 
解.
()
()
()
()
f
x
f x f
x
f
x
''
'
=
=
,
()
()
()
()
3!
3!
f
x
f
x f
x
f
x
'''
'
=
=
, 
()()
()
4!
f
x
f
x
=
,由归纳法,
()()
()
!
n
n
f
x
n f
x
+
=
,故
()()
!
n
f
n
=
. 
\textcolor{lyred}{例}.利用变换
cos
x
t
=
化简微分方程(
)
x
y
xy
y
''
'
-
-
+
=
. 
解.
sin
dy
dy dt
dy
dx
dt dx
dt
t
-
=
=
\cdots 
,
cos
sin
sin
sin
d y
d
dy
dt
d y
dy
t
dx
dt
dx
dx
dt
t
dt
t
t


-
-


=
=
\cdots 
+
\cdots 
\cdots 
=








cos
sin
sin
d y
dy
t
dt
t
dt
t
\cdots 
-
\cdots 
,代入方程,得
d y
y
dt
+
=
. 
\textcolor{lyred}{补充练习 }
1.设
()
()
f x
f
x
e
-
'
=
,
()
f
=,求
()()
n
f
. 
解.
()
()
()
()
f x
f x
f
x
e
f
x
e
-
-
''
'
=-
\cdots 
=-
,
()
()
()
()
2!
2!
f x
f x
f
x
e
f
x
e
-
-
'''
'
=
\cdots 
=
, 
()()
()
3!
f x
f
x
e
-
=-
,
()()
(
)
(
)
()
1 !
n
n
n f x
f
x
n
e
-
-\cdots 
=-
-
,故
()()
(
)
(
)
1 !
n
n
n
f
n
e
-
-
=-
-
. 
2.利用变换
t
x
e
=
化简微分方程
x y
xy
y
''
'
+
-
=
. 
解.
t
dy
dy dt
dy e
dx
dt dx
dt
-
=
=
,
t
t
t
t
t
d y
d
dy
dt
d y
dy
d y
dy
e
e
e
e
e
dx
dt dt
dx
dt
dt
dt
dt
-
-
-
-
-






=
=
-
=
-












, 
代入方程,得
d y
y
dt
-
=
. 
\vspace{6pt}
{\large \bfseries \textcolor{lyred}{第2.4 节 隐函数及由参数方程所确定的函数的导数 相关变化率}}

\vspace{4pt}
{\bfseries \textcolor{lyred}{一.隐函数的导数}}

\textcolor{lyred}{例}.设
()
y
y x
=
是方程
ye
xy
e
+
-
=
确定的隐函数,求dy
dx . 
解.方程两边对x 求导,得
y dy
dy
e
y
x
dx
dx
+
+
=
,故
y
dy
y
dx
x
e
=-
+
. 
\textcolor{lyred}{例}.设
()
y
y x
=
是方程
y
y
x
x
+
-
-
=
确定的隐函数,求
x
dy
dx
=
. 
解.
dy
dy
dy
x
y
x
dx
dx
dx
y
+
+
--
=
\Rightarrow 
=
+
,故
(
)
0,0
dy
dx
=
. 
\textcolor{lyred}{例}.求椭圆
x
y
+
=在
2,






处的切线斜率. 
解. 2
x
y dy
dy
x
dx
dx
y
+
\cdots 
=
\Rightarrow 
=-
\cdots 
,代入
x =
,
y =
,得
k =-
. 
\textcolor{lyred}{例(对数求导法)}.求
(
)
sin
x
y
x
x
=
>
的导数. 
解.ln
sin
ln
y
x
x
=
\cdots 
,两边对x 求导,得1
cos
ln
sin
y
x
x
x
y
x
'=
\cdots 
+
\cdots 
,故 
sin
sin
sin
cos
ln
cos
ln
x
x
x
y
y
x
x
x
x
x
x
x




'=
\cdots 
+
=
\cdots 
+








. 
\textcolor{lyred}{例}.设
()
y
y x
=
是方程
y
x
x
y
+
=确定的隐函数,求y'. 
解.
ln
ln
ln
ln
ln
ln
y
x
y
x
x
y
y
x
x
y
y
y
x
y
e
e
e
y
x
e
y
x
x
y
'




'
+
=\Leftrightarrow 
+
=\Rightarrow 
+
+
+
\cdots 
=








,故 
ln
ln
x
y
y
x
y
y
yx
y
x
x
xy
-
-
+
'=-
+
. 
\textcolor{lyred}{例}.设
(
)(
)
(
)(
)
x
x
y
x
x
-
-
=
-
-
,求y'. 
解.
(
)
ln
ln
ln
ln
ln
y
x
x
x
x
=
-
+
-
-
-
-
-
,两边对x 求导,得 
y
y
x
x
x
x
'


=
+
-
-


-
-
-
-


,故
y
y
x
x
x
x


'=
+
-
-


-
-
-
-


. 
\textcolor{lyred}{例}.设
(
)
(
)
+
-
+
=
x
x
x
y
,求y'. 
解.
ln
ln
4ln 3
5ln
y
x
x
x
=
+
+
-
-
+
,
y
y
x
x
x
'
-
=
\cdots 
+
-
+
-
+,故 
(
)
(
)
2 3
x
x
y
y
x
x
x
x
x
x
x
+
-




'=
+
-
=
+
-




+
-
+
+
-
+




+
. 
\textcolor{lyred}{例}.设
()
y
y x
=
是方程2
sin
y
x
y
-
+
=
确定的隐函数,求y''. 
解.
cos
dy
dx
y
=
+
,
(
)
(
)
(
)
(
)
2 2
cos
2sin
4sin
cos
cos
cos
y
d y
y y
y
dx
y
y
y
'
-
+
'
\cdots 
=
=
=
+
+
+
. 
\textcolor{lyred}{例}.设
()
y
y x
=
是方程
y
y
xe
=-
确定的隐函数,求y''. 
解.
y
y
e
y
xe
'=-+
,
(
)
(
)
(
)
(
)
y
y
y
y
y
y
y
y
y
e
y
xe
e
e
xe
y
e
xe
y
xe
xe
'
'
\cdots 
+
-
+
\cdots 
+
''=-
=
+
+
. 
\vspace{4pt}
{\bfseries \textcolor{lyred}{二.由参数方程所确定的函数的导数}}

\textcolor{lyred}{公式}.设
()
()
x
x t
y
y t
=

=

,若
()
x t
'
\neq 
,则
dy
dy
dt
dx
dx
dt
=
. 
\textcolor{lyred}{例}.设
()
y
y x
=
是
cos
x
t
y
t
=+

=

确定的函数,求dy
dx . 
解.
sin
dy
t
dx
t
-
=
. 
\textcolor{lyred}{例}.求阿基米德螺线
(
)
a
a
\rho 
\theta 
=
>
在
\pi 
\theta =
处的切线斜率. 
解.由
cos
sin
x
y
\rho 
\theta 
\rho 
\theta 
=


=

,得
cos
sin
x
a
y
a
\theta 
\theta 
\theta 
\theta 
=


=

,
sin
cos
cos
sin
dy
a
a
dx
a
a
\theta 
\theta 
\theta 
\theta 
\theta 
\theta 
+
=
-
,故
k
\pi 
\pi 
+
=
-
. 
\textcolor{lyred}{例}.设
()
y
y x
=
是
sin
1 cos
x
t
t
y
t
=-


=-

确定的函数,求
dx
y
d
. 
解.
sin
1 cos
dy
t
dx
t
=-
,
(
)
(
)
(
)
cos
1 cos
sin
1 cos
1 cos
t
t
t
d y
dt
dx
dx
t
t
-
-
-
=
\cdots 
=
-
-
. 
\textcolor{lyred}{例}.设
()
()
()
x
f
t
y
tf
t
f t
'
=

'
=
-

,其中
()
f t 二阶可导,
()
f
t
''
\neq 
,求
d y
dx . 
解.
()
()
()
()
f
t
tf
t
f
t
dy
t
dx
f
t
'
''
'
+
-
=
=
''
,
()
d y
dt
dx
dx
f
t
=\cdots 
=
''
. 
\vspace{4pt}
{\bfseries \textcolor{lyred}{三.相关变化率问题}}

\textcolor{lyred}{例}.设气球从离观察员500m 处垂直上升,当气球高度为500m 时,测得上升速度 
为140m / min ,问此时观察员视线仰角的增加率是多少? 
解.设气球上升t 分钟后,高度为h 米,观察员视线的仰角为\alpha ,则tan
h
\alpha =
, 
对t 求导,得
sec
d
dh
dt
dt
\alpha 
\alpha \cdots 
=
\cdots 
,代入
dh
dt =
,
\alpha =
\circ ,得
d
dt
\alpha 
\cdots 
=
\cdots 
,故 
0.14
d
dt
\alpha =
(弧度/分). 
\textcolor{lyred}{例}.设有深8m 顶部直径8m 的正圆锥形容器,现以
4m / min 的速度往其中注水, 
求水深为5m 时,容器内水面上升的速度. 
解.当水深h 时,水面半径
r
h
=
,面积
S
h
\pi 
=
,故此时容器中水的体积为 
V
Sh
h
\pi 
=
=
,两边对t 求导,得
dV
dh
dh
dV
h
dt
dt
dt
h
dt
\pi 
\pi 
=
\cdots 
\cdots 
\Rightarrow 
=
\cdots 
,代入 
h =
,
dV
dt =
,得
(
)
m / min
dh
dt
\pi 
=
. 
\textcolor{lyred}{补充练习 }
1.设
()
y
y x
=
是
sin
xe
t
t
t
y
y
\pi 

=
+
+

-
+
=

确定的函数,求
t
dy
dx
=
. 
解.
(
)
(
)(
)
sin
sin
2 1
cos
sin
cos
cos
x
x
x
dx
dx
e
t
e
t
dy
e
y
dt
dt
dy
y
dy
dy
dx
t
t
y
y
t
y
dt
t
y
dt
dt
-


=
+
=
+



\Rightarrow 
\Rightarrow 
=


+
-


=
+
-
=

-



,代入 
t =
,
x =
,
y
\pi 
=
,得
t
dy
dx
=
=
. 
\vspace{6pt}
{\large \bfseries \textcolor{lyred}{第2.5 节 函数的微分}}

\vspace{4pt}
{\bfseries \textcolor{lyred}{一.微分的定义}}

\textcolor{lyred}{例}.边长x 的正方形,面积
A
x
=
,若x 增加
x
\Delta ,则相应的面积增量 
(
)
(
)
(
)
A
x
x
x
x x
x
x x
o
x
\Delta 
\Delta 
\Delta 
\Delta 
\Delta 
\Delta 
=
+
-
=
+
=
+
,其中2x x
\Delta 是x
\Delta 的\textcolor{lyred}{线性函数}, 
并且是A
\Delta 的\textcolor{lyred}{主要部分},因为两者之间的误差是x
\Delta 的高阶无穷小,记
dA
x x
=
\Delta , 
称为
2x 在x 处的\textcolor{lyred}{微分}. 
\textcolor{lyred}{定义}.设
()
y
f x
=
在某个
()
U x 内有定义,若
(
)
()
y
f x
x
f x
\Delta 
\Delta 
=
+
-
可以表示为: 
()
(
)
y
x
x
o
x
\Delta 

\Delta 
\Delta 
=
+
,即x
\Delta 的线性函数与高阶无穷小之和,其中的
()
A x 与
x
\Delta  
无关,则称
()
f x 在x 处\textcolor{lyred}{可微},记
()
dy
x
x

\Delta 
=
,称为
()
f x 在x 处相对于自变量的 
增量x
\Delta 的\textcolor{lyred}{微分},也可记为
()
df x . 
\vspace{4pt}
{\bfseries \textcolor{lyred}{二.可微的条件}}

\textcolor{lyred}{定理}.
()
y
f x
=
在x 处可微
()
y
f x
\Leftrightarrow 
=
在x 处可导. 
\textcolor{lyred}{注}.设
()
y
f x
=
在x 处可微,则
()
()
dy
df x
f
x
x
\Delta 
'
=
=
. 
\textcolor{lyred}{注(有限增量公式)}.
()
y
f
x
x
x
\Delta 
\Delta 
\alpha \Delta 
'
=
+
\cdots 
,其中
lim
x
\Delta 
\alpha 
\to 
=
. 
\textcolor{lyred}{注}.若
()
f
x
'
\neq 
,则
()
lim
lim
x
x
y
y
dy
f
x
x
\Delta \to 
\Delta \to 
\Delta 
\Delta 
=
=
'
\Delta 
,即当
x
\Delta \to 
时,
~
y
dy
\Delta 
,于是 
(
)
y
dy
o dy
\Delta 
=
+
,称dy 为y
\Delta 的\textcolor{lyred}{线性主部}. 
\textcolor{lyred}{定义}.若
()
y
f x
=
在区间I 上处处可微,则称它为I 上的\textcolor{lyred}{可微函数};此时,记 
()
()
dy
df x
f
x
x
\Delta 
'
=
=
,称为\textcolor{lyred}{函数的微分}. 
\textcolor{lyred}{注}.
()
dx
x
x
x
'
=
\Delta =\Delta ,即自变量的微分等于它的增量,故
()
dy
f
x dx
'
=
. 
\textcolor{lyred}{例}.求
y
x
=
在
x =
处当
0.1
x
\Delta =
时的微分. 
解.
(
)
dy
d x
x
x
\Delta 
=
=
,故
3,
0.1
108 0.1
10.8
x
x
dy
\Delta 
=
=
=
\cdots 
=
. 
\vspace{4pt}
{\bfseries \textcolor{lyred}{三.微分的几何意义}}

曲线
()
y
f x
=
上点(
)
,x y 处对应于横坐标增量x
\Delta 的纵坐标增量为
y
\Delta ,而曲线在 
(
)
,x y 处的切线上相应的纵坐标增量为dy ,当x
\Delta 很小时,两者近似相等. 
\vspace{4pt}
{\bfseries \textcolor{lyred}{四.微分运算法则}}

\textcolor{lyred}{法则1}. (
)
d
u
v
du
dv
\alpha 
\beta 
\alpha 
\beta 
\pm 
=
\pm 
, (
)
d uv
vdu
udv
=
+
,
u
vdu
udv
d
v
v
-

=




.\textcolor{lyred}{ }
\textcolor{lyred}{法则2}.设
()
y
f u
=
,
()
u
x
\varphi 
=
均可微,则
()
()
()
dy
f
g x
dg x
f
u du
'
'
=
=




.\textcolor{lyred}{ }
\textcolor{lyred}{注}.设
()
y
f u
=
,则无论u 是自变量,还是中间变量,微分的表达式
()
dy
f
u du
'
=
总是正确的,这个性质称为\textcolor{lyred}{一阶微分的形式不变性}. 
\textcolor{lyred}{例}.
(
)
cos
cos
sin
cos
2 2
x
x
d
x
x
d
x
d
x
dx
x
x
+
+
+=
+
+=
+
=
+
+
. 
\textcolor{lyred}{例}.
(
)
(
)
(
)
ln 1
x
x
x
x
x
x
x
e
xe
d
e
d
e
d x
dx
e
e
e
+
=
+
=
=
+
+
+
. 
\textcolor{lyred}{例}. (
)
(
)
1 3
1 3
1 3
1 3
cos
cos
cos
3cos
sin
x
x
x
x
d e
x
x de
e
d
x
e
x
x dx
-
-
-
-
=
\cdots 
+
\cdots 
=-
+
. 
\textcolor{lyred}{例}.设
()
y
y x
=
是
(
)
sin
cos
y
x
x
y
-
-
=
确定的隐函数,求dy . 
解. (
)
(
)
(
)(
)
sin
cos
sin
sin
sin
d y
x
d
x
y
xdy
yd
x
x
y d x
y
-
-
=
\Rightarrow 
+
+
-
-
=
,解得 
(
)
(
)
cos
sin
sin
sin
y
x
x
y
dy
dx
x
y
x
+
-
=
-
-
. 
\vspace{4pt}
{\bfseries \textcolor{lyred}{五.函数的近似计算}}

\textcolor{lyred}{公式}.设
(
)
f
x
'
\neq 
,则
()
(
)
(
)(
)
f x
f x
f
x
x
x
'
\approx 
+
-
,误差为(
)
o x
x
-
. 
\textcolor{lyred}{注}.特别地,若
()
f '
\neq 
,则当
x <<时,
()
()
()
f x
f
f
x
'
\approx 
+
. 
\textcolor{lyred}{例}.求
0.97 和1.05 的近似值. 
解.
()
f x
x
=
,
x =,
0.97
x =
,
0.03
x
\Delta =-
,
()
x
f
x
=
'
=
=
,于是 
(
)
()
()
0.97
0.97
0.03
0.985
f
f
f
x
\Delta 
'
=
\approx 
+
=-
\cdots 
=
,查表:0.9849 , 
(
)
()
()
1.05
1.05
0.05
1.025
f
f
f
x
'
=
\approx 
+
\Delta =+
\cdots 
=
,查表:1.0247 . 
\textcolor{lyred}{例}.在半径为1cm 的球上镀一层厚度为0.01cm 的铜,估计每个球需用铜多少
cm . 
解.
V
R
\pi 
=
,
R =,
0.01
R
\Delta 
=
,故镀层体积
(
)
(
)
V
V R
R
V R
\Delta 
=
+\Delta 
-
\approx  
(
)
4 3.14 1 0.01
0.13cm
V
R
R
'
\Delta 
=
\cdots 
\cdots 
\cdots 
\approx 
. 
\textcolor{lyred}{补充练习 }
1.函数
()
f x 在
x =
处可导的充要条件是______. 
(A)
()
f x 在
x =
处连续;        (B)
()
()
()
f x
f
Ax
o x
-
=
+
; 
(C)
()
f-'
与
()
f+'
均存在;       (D)
()
lim
x
f
x
\to 
'
存在. 
解.选(B). 